
\documentclass[12pt, a4paper]{article}
\usepackage[utf8]{inputenc}
\usepackage{amsmath}
\usepackage{amssymb} % For \Box
\usepackage{graphicx}
\usepackage[a4paper, margin=1in]{geometry}
\usepackage{sectsty}
\usepackage{hyperref}
\usepackage{titlesec}
\usepackage{parskip} 
\usepackage[dvipsnames,svgnames,x11names]{xcolor}

% Customizing titles
\titleformat{\section}
  {\normalfont\fontsize{18}{22}\bfseries\color{MidnightBlue}}
  {\thesection}
  {1em}
  {}
\titleformat{\subsection}
  {\normalfont\fontsize{14}{18}\bfseries\color{black!80}}
  {\thesubsection}
  {1em}
  {}
\titleformat{\section*}
  {\normalfont\fontsize{18}{22}\bfseries\color{MidnightBlue}}
  {}
  {0em}
  {}
\titleformat{\subsection*}
  {\normalfont\fontsize{14}{18}\bfseries\color{black!80}}
  {}
  {0em}
  {}

\hypersetup{
    colorlinks=true,
    linkcolor=MidnightBlue,
    urlcolor=DarkCyan,
}

\usepackage{fancyhdr}
\pagestyle{fancy}
\fancyhf{}
\fancyhead[C]{Chemistry SME Agent}
\fancyfoot[C]{Page \thepage}
\renewcommand{\headrulewidth}{0.4pt}
\renewcommand{\footrulewidth}{0.4pt}

\title{periodic table}
\author{Generated by AI Educational Assistant}
\date{}

\begin{document}
\maketitle
\thispagestyle{fancy}

\section*{Report on: periodic table}

\subsection*{Introduction}
The periodic table is a fundamental framework in chemistry, representing the ordered arrangement of all known chemical elements. This systematic organization is based on increasing atomic number, which corresponds to the total number of protons in an atom's nucleus. More than just a list, the periodic table embodies the "periodic law," a profound principle describing recurring patterns in which elements sharing the same column, or group, exhibit comparable chemical and physical properties. Its development, initially pioneered by Dmitry I. Mendeleyev in the mid-19th century, has had an immeasurable impact on the advancement of chemical understanding and continues to be an indispensable tool for students and scientists alike.

\subsection*{Key Findings}
The modern periodic table is defined by the arrangement of chemical elements in order of increasing atomic number (Z). The atomic number is equivalent to the total number of protons within an atom's nucleus, a concept that was not fully understood until the second decade of the 20th century. This understanding clarified the basis of the periodic law, which posits that when elements are grouped by atomic number, a recurring pattern emerges where elements in the same column (group) possess similar chemical characteristics.

Historically, the first periodic table was developed by the Russian chemist Dmitri Mendeleev. His original arrangement, however, grouped elements according to their rising atomic weight. Significant advancements in the early 20th century led to a better understanding of the electrical structure of atoms and molecules, revealing that atomic number, not atomic weight, is the true determinant of an element's position and properties within the periodic system. This shift was crucial for the development of the contemporary periodic table, which is now recognized as a vital framework for organizing and understanding chemical elements.

Currently, the periodic table contains a total of 118 elements. Its structure and content are overseen by the International Union of Pure and Applied Chemistry (IUPAC). The table serves not only as a means to identify elements easily but also to reveal repeating trends in their chemical and physical properties, which are essential for studying atomic structure, element classification, and periodic trends in chemistry.

\begin{itemize}
 \item \textbf{First and Last Elements:} Hydrogen (H) is the first element, and Oganesson (Og) is currently the last element listed.
 \item \textbf{Most Prevalent Elements:} Hydrogen (H) is the most prevalent element in the universe, accounting for approximately 75\% of the observable cosmos. Helium (He) is the second most prevalent, with hydrogen and helium together making up over 99\% of the observable universe.
 \item \textbf{First Artificial Element:} Technetium (Tc) holds the distinction of being the first element to be created artificially.
 \item \textbf{Governing Principle:} The arrangement is based on increasing atomic number, which is the number of protons in an atom.
 \item \textbf{Periodic Law:} Elements in the same column (group) exhibit comparable properties due to recurring patterns in their atomic structure.
\end{itemize}

\subsection*{Conclusion}
In summary, the periodic table is an indispensable scientific tool, providing an ordered arrangement of all chemical elements based on their increasing atomic number. This structure underpins the "periodic law," which elucidates the recurring patterns of chemical and physical properties among elements within the same group. While initially conceived by Dmitri Mendeleev based on atomic weight, the modern periodic table's reliance on atomic number, understood in the early 20th century, solidified its scientific accuracy and predictive power. With 118 known elements and its governance by IUPAC, the periodic table remains a crucial framework for comprehending atomic structure, classifying elements, and predicting chemical behavior, making it fundamental to the study of chemistry across all levels.

\end{document}
